\section{Amplitudes}
\label{app:Amplitudes}

In this Appendix, we report the amplitudes that we employed throughout the main text, expressed in terms of spinor-helicity variables.

\subsection{3-point tree amplitudes}
Here, the analytically continued 3-point tree amplitudes in the holomorphic ($\hol$) and antiholomorphic ($\ahol$) configurations belonging to the different sectors of the Lagrangian are displayed on the left and on the right, respectively.
They are completely constrained, up to an overall factor (the coupling constant), by locality, Poincaré invariance, and dimensional analysis~\cite{Benincasa:2007xk}. Indeed, in full generality they read as follows
\begin{align}
    \mathcal M_3^\hol(1^{h_1},2^{h_2},3^{h_3}) = g_\hol \agl{1}{2}^{a_3}\agl{2}{3}^{a_1}\agl{3}{1}^{a_2}\,, &&
    \mathcal M_3^\ahol(1^{h_1},2^{h_2},3^{h_3}) = g_\ahol \sqr{1}{2}^{\bar a_3}\sqr{2}{3}^{\bar a_1}\sqr{3}{1}^{\bar a_2}\,,
\end{align}
with $\bar a_i=-a_i$,
\begin{align}
    a_1= h_1-h_2-h_3\,,&&
    a_2= h_2-h_3-h_1\,, &&
    a_3= h_3-h_1-h_2\,,
\end{align}
and the mass dimensions of the coupling constants only depend on the helicities:
\begin{align}
    [g_\hol]=1+h_1+h_2+h_3\,, && 
    [g_\ahol]=1-h_1-h_2-h_3\,.
\end{align}
Locality implies $[g_\hol],[g_\ahol]<1$, therefore we can infer that the holomorphic (antiholomorphic) configuration is the consistent one if $h_1+h_2+h_3<0$ ($h_1+h_2+h_3>0$).
The case where $h_1+h_2+h_3=0$ is trivial, as it can only correspond to a cubic scalar interaction, where $h_1=h_2=h_3=0$.

\paragraph{$\mathcal L_{\text{LO}}$.}
The lowest order Lagrangian
\begin{equation}
    \mathcal L_{\text{LO}} = -\frac{1}{4}G_{\mu\nu}^a G^{a,\mu\nu}-\frac{1}{4}F_{\mu\nu}F^{\mu\nu}+i\bar f_i \gamma^\mu D_\mu f_i - y_i h \bar f_i f_i
\end{equation}
generates
\begin{align}
    \mathcal M(1^-_{f_i},2^+_{\bar f_j},3^-_\gamma) &= \sqrt 2 e Q_f \delta^{ij} \frac{\agl{1}{3}^2}{\agl{1}{2}}\,, &
     \mathcal M(1^-_{f_i},2^+_{\bar f_j},3^+_\gamma) &= \sqrt 2 e Q_f \delta^{ij}\frac{\sqr{3}{2}^2}{\sqr{1}{2}}\,, \\
     \mathcal M(1^-_{f_i^I},2^+_{\bar f_j^J},3^-_{g^a}) &= \sqrt 2 g_s c_f T^a_{IJ}\delta^{ij} \frac{\agl{1}{3}^2}{\agl{1}{2}}\,, &
     \mathcal M(1^-_{f_i^I},2^+_{\bar f_j^J},3^+_{g^a}) &= \sqrt 2 g_s c_f T^a_{IJ} \delta^{ij}\frac{\sqr{3}{2}^2}{\sqr{1}{2}}\,, \\
     \mathcal M(1^-_{g^a},2^-_{g^b},3^+_{g^c}) &= i\sqrt 2 g_s f^{abc}\frac{\agl{1}{2}^3}{\agl{1}{3}\agl{3}{2}}\,,  &
     \mathcal M(1^+_{g^a},2^+_{g^b},3^-_{g^c}) &=-i\sqrt 2 g_s f^{abc}\frac{\sqr{1}{2}^3}{\sqr{1}{3}\sqr{3}{2}}\,,  \\
     \mathcal M(1^-_{f_i},2^-_{\bar f_j},3_h) &= - y_i\delta^{ij}\agl{1}{2}\,, &
     \mathcal M(1^+_{f_i},2^+_{\bar f_j},3_h) &= - y_i \delta^{ij} \sqr{1}{2} \,.
\end{align}

\paragraph{$\mathcal L_{\phi}$.}
The ALP effective Lagrangian
\begin{equation}
    \mathcal{L}_\phi = 
    \frac{\tilde{\mathcal{C}}_\gamma}{\Lambda}\,\phi\, F\tilde F
    + \frac{\tilde{\mathcal{C}}_g}{\Lambda}\,\phi\, G\tilde G
    + \mathcal{Y}^{ij}_P\,\phi\, \bar{f}_i i\gamma_5 f_j 
     + \frac{\mathcal{C}_\gamma}{\Lambda}\,\phi\, FF
    + \frac{\mathcal{C}_g}{\Lambda}\,\phi\, GG
    + \mathcal{Y}^{ij}_S\,\phi\, \bar{f}_i f_j
\end{equation}
generates
\begin{align}
    \mathcal M(1^-_\gamma,2^-_\gamma,3_\phi) &= -\frac{2}{\Lambda}(\mathcal C_\gamma + i \tilde{\mathcal{C}}_\gamma)\agl{1}{2}^2\,, &
    \mathcal M(1^+_\gamma,2^+_\gamma,3_\phi) &= -\frac{2}{\Lambda}(\mathcal C_\gamma - i \tilde{\mathcal{C}}_\gamma)\sqr{1}{2}^2\,, \\
    \mathcal M(1^-_{g^a},2^-_{g^b},3_\phi) &= -\frac{2}{\Lambda}(\mathcal C_g + i \tilde{\mathcal{C}}_g)\delta^{ab}\agl{1}{2}^2\,, &
    \mathcal M(1^+_{g^a},2^+_{g^b},3_\phi) &= -\frac{2}{\Lambda}(\mathcal C_g - i \tilde{\mathcal{C}}_g)\delta^{ab}\sqr{1}{2}^2\,, \\
    \mathcal M(1^-_{f_i},2^-_{\bar f_j},3_\phi)&=(\mathcal Y_S^{ij}-i\mathcal Y_P^{ij})\agl{1}{2}\,, &
    \mathcal M(1^+_{f_i},2^+_{\bar f_j},3_\phi)&=(\mathcal Y_S^{ij}+i\mathcal Y_P^{ij})\sqr{1}{2} \,.
\end{align}

\paragraph{$\mathcal L^{(5)}$.}
The relevant dimension-$5$ Lagrangian invariant under $SU(3)_c\times U(1)_{\text{em}}$ and built of SM particles consists of dipole operators
\begin{equation}
    \mathcal L^{(5)} = \frac{c_{\text{M}}^{ij}}{\Lambda}\bar f_i \sigma^{\mu\nu} f_j F_{\mu\nu} + \frac{c_{\text{E}}^{ij}}{\Lambda}\bar f_i \sigma^{\mu\nu}i\gamma_5 f_j F_{\mu\nu} + \frac{c_{\text{CM}}^{ij}}{\Lambda}\bar f_i \sigma^{\mu\nu}T^a f_j G^a_{\mu\nu} + \frac{c_{\text{CE}}^{ij}}{\Lambda}\bar f_i \sigma^{\mu\nu}i\gamma_5 T^a f_j G^a_{\mu\nu}
\end{equation}
and generates
\begin{align}
    \mathcal M(1^-_{f_i},2^-_{\bar f_j},3^-_\gamma) &= \frac{2\sqrt 2}{\Lambda}C_\gamma^{ij}\agl{1}{3}\agl{2}{3}\,, &
    \mathcal M(1^+_{f_i},2^+_{\bar f_j},3^+_\gamma)&=- \frac{2\sqrt 2}{\Lambda}(C_\gamma^{ij})^*\sqr{1}{3}\sqr{2}{3}\,,  \\
    \mathcal M(1^-_{f_i^I},2^-_{\bar f_j^J},3^-_{g^a}) &= \frac{2\sqrt 2}{\Lambda}C_g^{ij}T^a_{IJ}\agl{1}{3}\agl{2}{3}\,, &
    \mathcal M(1^+_{f_i^I},2^+_{\bar f_j^J},3^+_{g^a})&= - \frac{2\sqrt 2}{\Lambda}(C_g^{ij})^*T^a_{IJ}\sqr{1}{3}\sqr{2}{3}\,,
\end{align}
with
\begin{align}
    C_\gamma^{ij}=c_{\text{M}}^{ij}-i\, c_{\text{E}}^{ij}\,, && C_g^{ij}=c_{\text{CM}}^{ij}-i\,c_{\text{CE}}^{ij}\,.
\end{align}

\paragraph{$\mathcal L^{(6)}$.}
The relevant dimension-$6$ Lagrangian invariant under $SU(3)_c\times U(1)_{\text{em}}$ and built of SM particles only consists of
\begin{equation}
    \mathcal L^{(6)} = \frac{D_G}{3\Lambda^2} f^{abc} G_\mu^{a, \nu} G_{\nu}^{b, \rho} G_{\rho}^{c, \mu} + \frac{d_G}{3\Lambda^2} f^{abc} G_\mu^{a, \nu} G_{\nu}^{b, \rho} \tilde G_{\rho}^{c, \mu}
\end{equation}
and generates
\begin{align}
    \mathcal M(1^-_{g^a},2^-_{g^b},3^-_{g^c}) = i\frac{\sqrt 2}{\Lambda^2}C_Gf^{abc}\agl{1}{2}\agl{2}{3}\agl{1}{3}\,, &&
     \mathcal M(1^+_{g^a},2^+_{g^b},3^+_{g^c}) = i\frac{\sqrt 2}{\Lambda^2}C_G^* f^{abc}\sqr{2}{1}\sqr{3}{2}\sqr{3}{1}\,,
\end{align}
with
\begin{equation}
    C_G= D_G + i\, d_G\,.
\end{equation}

\subsection{4-point tree amplitudes}
Here, the 4-point tree amplitudes needed for the calculations are displayed.
With the symbol $*$ we denote the region in the space of the couplings of the theory where only the gauge couplings $e$ and $g_s$ are different from zero.
%
\begin{align}
     \mathcal M|_*(1^-_{f_i},2^+_{\bar f_j},3^-_\gamma,4^+_\gamma) &= -2e^2Q_f^2\delta^{ij}\frac{\agl{1}{3}\sqr{4}{2}}{\agl{1}{4}\sqr{3}{1}}\,, \\
    \mathcal M|_*(1^+_{f_i},2^-_{\bar f_j},3^-_\gamma,4^+_\gamma) &= -2e^2Q_f^2\delta^{ij}\frac{\agl{2}{3}\sqr{4}{1}}{\agl{1}{4}\sqr{3}{1}}\,, \\
    \mathcal M|_*(1^-_{f_i^I},2^+_{\bar f_j^J},3^-_{g^a},4^+_{g^b}) &= -2g_s^2c_f^2T^a_{IK}T^b_{KJ}\delta^{ij}\frac{\agl{1}{3}\sqr{4}{2}}{\agl{1}{4}\sqr{3}{1}}\,, \\
    \mathcal M|_*(1^+_{f_i^I},2^-_{\bar f_j^J},3^-_{g^a},4^+_{g^b}) &= -2g_s^2c_f^2T^b_{IK}T^a_{KJ}\delta^{ij}\frac{\agl{2}{3}\sqr{4}{1}}{\agl{1}{4}\sqr{3}{1}}\,, \\
    \mathcal M|_*(1^-_{g^a},2^-_{g^b},3^+_{g^c},4^+_{g^d}) &= -2g_s^2\agl{1}{2}^4\bigg(\frac{f^{abe}f^{cde}}{\agl{1}{2}\agl{2}{3}\agl{3}{4}\agl{4}{1}}+\frac{f^{ace}f^{bde}}{\agl{1}{3}\agl{3}{2}\agl{2}{4}\agl{4}{1}}\bigg)\,,\\
    \mathcal M|_*(1^-_{f_i^I},2^-_{\bar f_j^J},3^+_{f_k^K},4^+_{\bar f_l^L}) &= -2 (e^2 Q_f^2 \delta_{IL}\delta_{KJ}+g_s^2c_f^2T^a_{IL}T^a_{KJ}) \delta^{il}\delta^{jk} \frac{\agl{1}{2}\sqr{4}{3}}{\agl{1}{4}\sqr{4}{1}}\,, \\
    \mathcal M|_*(1^-_{f_i^I},2^+_{\bar f_j^J},3^-_{f_k^K},4^+_{\bar f_l^L}) &= +2(e^2Q_f^2\delta_{IJ}\delta_{KL}+g_s^2c_f^2T^a_{IJ}T^a_{KL})\delta^{ij}\delta^{kl} \frac{\agl{1}{3}\sqr{4}{2}}{\agl{1}{2}\sqr{2}{1}}\nonumber \\
    &\quad - 2(e^2Q_f^2\delta_{IL}\delta_{KJ}+g_s^2c_f^2T^a_{IL}T^a_{KJ}) \delta^{il}\delta^{jk} \frac{\agl{1}{3}\sqr{4}{2}}{\agl{1}{4}\sqr{4}{1}}\,, \\
    \mathcal M|_*(1^-_{f_i^I},2^+_{\bar f_j^J},3^+_{f_k^K},4^-_{\bar f_l^L}) &= +2(e^2Q_f^2\delta_{IJ}\delta_{KL}+g_s^2c_f^2T^a_{IJ}T^a_{KL})\delta^{ij}\delta^{kl} \frac{\agl{1}{4}\sqr{3}{2}}{\agl{1}{2}\sqr{2}{1}}\,,\\
    \mathcal M(1^-_{f_i},2^+_\gamma,3^-_{\bar f_j},4_h) &= -\sqrt{2} y_i \delta^{ij} e Q_f \frac{\agl{1}{3}^2}{\agl{1}{2}\agl{2}{3}}\,, \\
    \mathcal M(1^-_{f_i^I},2^+_{g^a},3^-_{\bar f_j^J},4_h) &= -\sqrt{2} y_i \delta^{ij} g_s c_f T^a_{IJ} \frac{\agl{1}{3}^2}{\agl{1}{2}\agl{2}{3}} \,,\\
    \mathcal M(1^-_{f_i},2^-_\gamma,3^+_{\bar f_j},4_\phi) &= \frac{2\sqrt 2}{\Lambda} e Q_f \delta^{ij}(\mathcal C_\gamma + i\tilde{\mathcal C}_\gamma)\frac{\agl{1}{2}^2}{\agl{1}{3}}\,,\\
    \mathcal M(1^-_{f_i^I},2^-_{g^a},3^+_{\bar f_j^J},4_\phi) &= \frac{2\sqrt 2}{\Lambda} g_s c_f T^a_{IJ} \delta^{ij}(\mathcal C_g + i\tilde{\mathcal C}_g)\frac{\agl{1}{2}^2}{\agl{1}{3}}\,,\\
    \mathcal M(1^-_{f_i},2^-_{\bar f_j},3^+_\gamma,4_\phi) &= -\sqrt 2 e Q_f (\mathcal Y_S^{ij}-i\mathcal Y_P^{ij})\frac{\agl{1}{2}^2}{\agl{1}{3}\agl{2}{3}} \,,\\
    \mathcal M(1^-_{f_i^I},2^-_{\bar f_j^J},3^+_{g^a},4_\phi) &= -\sqrt 2 g_s c_f T^a_{IJ} (\mathcal Y_S^{ij}-i\mathcal Y_P^{ij})\frac{\agl{1}{2}^2}{\agl{1}{3}\agl{2}{3}}\,,\\
    \mathcal M(1^-_{g^a},2^-_{g^b},3^+_{g^c},4_\phi) &= -i\frac{2\sqrt 2}{\Lambda}g_sf^{abc}(\mathcal C_g+i\tilde{\mathcal C}_g)\frac{\agl{1}{2}^3}{\agl{1}{3}\agl{2}{3}}\,.
\end{align}
